\documentclass[12pt]{article}

\usepackage{sbc-template}
\usepackage{graphicx,url}
\usepackage[utf8]{inputenc}
\usepackage[brazilian]{babel}
\usepackage{booktabs}
\usepackage{multirow}
\usepackage{amsmath}

\sloppy

\title{Construção e Avaliação Preliminar de um \textit{Dataset} para Análise da Complexidade Técnico-Linguística de Descrições do SUS}

\author{Lucas Matheus Santos da Silva\inst{1}, Pedro Ivo Leite\inst{1}}

\address{Programa de Pós-Graduação em Tecnologia da Informação (PPGTI)\\
  Instituto Federal da Paraíba (IFPB) -- João Pessoa, PB -- Brasil\\
  \email{lucas.silva@academico.ifpb.edu.br, pedro.leite@academico.ifpb.edu.br}
}

\begin{document}

\maketitle

\begin{abstract}
Official descriptions of diagnoses and procedures in the Brazilian Unified
Health System (SUS) contain specialized terminology, acronyms and
abbreviations that may hinder comprehension by non-specialists. This work
constructs a reproducible dataset for the exploratory analysis of
technical-linguistic complexity in ICD-10 and SIGTAP descriptions. The corpus
contains 17{,}182 records: 12{,}198 ICD-10 subcategories and 4{,}984 SIGTAP
procedures. We describe two pipeline variants: an \emph{executed} variant that
identifies medical jargon through Greek-Latin morphological signals, and a
\emph{proposed} variant that replaces the jargon proxy with DeCS-based
terminology matching once the licensed XML is released. Complexity bands are
interpreted as heuristic until compared with human judgments. Preliminary
descriptive analysis shows that source-level differences cannot be reduced to
text length alone.
\end{abstract}

\begin{resumo}
Descrições oficiais de diagnósticos e procedimentos do Sistema Único de Saúde
(SUS) contêm terminologia especializada, siglas e abreviações que podem
dificultar a compreensão por pessoas não especialistas. Este trabalho constrói
um \textit{dataset} reproduzível para análise exploratória da complexidade
técnico-linguística de descrições da CID-10 e do SIGTAP. O corpus possui
17.182 registros: 12.198 subcategorias CID-10 e 4.984 procedimentos SIGTAP.
Descrevemos duas variantes do \textit{pipeline}: uma \emph{executada}, que
identifica jargão médico por sinais morfológicos greco-latinos, e uma
\emph{proposta}, que substitui o \textit{proxy} de jargão por correspondência
terminológica com o DeCS quando o XML licenciado for liberado. As faixas de
complexidade são interpretadas como heurísticas até serem comparadas com
julgamentos humanos. A análise descritiva preliminar mostra que diferenças
entre fontes não podem ser reduzidas apenas ao comprimento do texto.
\end{resumo}

\section{Introdução}

A comunicação em saúde é fator determinante para a adesão a tratamentos e para
a segurança do paciente, e é reconhecidamente influenciada pela
\textit{health literacy} do indivíduo~\cite{nutbeam2008,sorensen2012}.
Diagnósticos e procedimentos do Sistema Único de Saúde (SUS) são registrados
em linguagem altamente técnica, com termos de origem greco-latina, siglas e
abreviações próprias da área médica. Para quem não é especialista ---
pacientes, estudantes e até profissionais de outras áreas da saúde ---
essa linguagem pode constituir uma barreira de compreensão. Revisões em
comunicação em saúde apontam que jargão e abreviações podem prejudicar a
interpretação de instruções por leitores não especialistas~\cite{maghroudi2021}.

Bases públicas do SUS, como a CID-10 e o SIGTAP, mantêm as descrições oficiais
desses diagnósticos e procedimentos. Apesar de amplamente utilizadas para fins
administrativos, essas descrições raramente são analisadas sob a ótica da
\textit{legibilidade}: o quão difícil é, para um leitor não especialista,
compreender aquele texto.

\textbf{Problema.} Este trabalho ataca o seguinte problema de Ciência de
Dados: \textit{construir um dataset que classifique, de modo heurístico,
diagnósticos e procedimentos médicos do SUS em níveis de complexidade
técnico-linguística --- baixa, média ou alta --- e identificar quais
atributos do texto mais influenciam essa classificação}. Trata-se de um
\textit{dataset} exploratório: a categoria é uma \textit{proxy} heurística,
não uma rotulagem validada por especialistas.

\textbf{Contribuições.} As principais contribuições são:
(i)~uma base, em português, construída a partir das tabelas oficiais do SUS;
(ii)~um \textit{pipeline} reproduzível de construção, enriquecimento e
formação de faixas heurísticas, descrito em duas variantes (executada e
proposta com DeCS); (iii)~uma análise descritiva inicial dos atributos por
fonte; e (iv)~um protocolo de validação com anotação humana.

Como \textit{recorte}, a variável de complexidade é criada por regras
heurísticas; assim, os resultados devem ser interpretados como uma análise
exploratória da complexidade textual, e não como medida definitiva da
compreensão real por parte dos pacientes.

Neste estudo, \textit{complexidade técnico-linguística} é uma definição
operacional baseada em terminologia de saúde, extensão, siglas, abreviações e
sinais morfológicos especializados. Ela não equivale à legibilidade em sentido
amplo, que também depende de sintaxe, contexto, apresentação e características
do leitor. As perguntas de pesquisa são: \textbf{QP1}: quais atributos
caracterizam descritivamente CID-10 e SIGTAP? \textbf{QP2}: como a terminologia
identificada pelo DeCS contribui para faixas exploratórias de complexidade?
\textbf{QP3}: qual é a concordância entre a regra heurística e uma amostra
estratificada anotada por humanos? Nesta entrega, \textbf{QP1} é parcialmente
respondida pelas estatísticas descritivas (Seção~\ref{sec:descritivos}); \textbf{QP2}
depende da liberação do XML do DeCS e \textbf{QP3} depende da rodada de anotação,
ambas tratadas como trabalho imediato.

\section{Trabalhos Relacionados}

A medição automática da dificuldade de leitura remonta às
\textit{fórmulas de legibilidade}. O índice de Flesch~\cite{flesch1948}
estima a legibilidade a partir do tamanho médio de frases e de palavras, e
foi adaptado ao português brasileiro por Martins et al.~\cite{martins1996},
com pesos próprios devido às diferenças entre os idiomas. A
fórmula de Dale-Chall~\cite{dalechall1948,chall1995} introduziu a noção
de \textit{palavras difíceis} --- aquelas fora de uma lista de palavras
familiares --- combinada ao tamanho das frases por meio de coeficientes
ponderados. Para o português, o Coh-Metrix-Port~\cite{scarton2010cohmetrix}
e, mais recentemente, o NILC-Metrix~\cite{leal2024nilcmetrix} reúnem dezenas
de métricas de complexidade textual, incluindo medidas lexicais, sintáticas e
de coesão. Trabalhos como o PorSimples~\cite{aluisio2010porsimples} aplicam
esses recursos a tarefas de simplificação textual em pt-BR.

No domínio médico, Zheng e Yu~\cite{zheng2017} mostraram empiricamente que
fórmulas clássicas de legibilidade subestimam a dificuldade percebida em
registros eletrônicos de saúde, sugerindo que o domínio exige métricas
específicas. O MedReadMe~\cite{jiang2024medreadme} avaliou centenas de
atributos e concluiu que a contagem de \textit{jargão} é o fator que mais
contribui para a dificuldade de textos médicos --- resultado que fundamenta
a priorização de termos técnicos neste trabalho. O estudo é inspiração
metodológica, não fórmula diretamente transferível: ele analisa sentenças em
inglês, enquanto este trabalho analisa descrições terminológicas em
português.

\textbf{Lacuna.} A literatura de legibilidade médica concentra-se no inglês e
em textos corridos. Não encontramos trabalho equivalente que classifique a
complexidade linguística das \textit{tabelas oficiais do SUS} em português.
É essa lacuna que o presente trabalho aborda.

\section{Metodologia}

O \textit{pipeline} de construção do \textit{dataset} segue seis etapas:
coleta, integração, limpeza, extração de atributos, classificação heurística
e análise. Foram utilizados exclusivamente dados públicos e secundários;
nenhum dado pessoal de paciente.

\paragraph{Variantes do pipeline.} Apresentamos duas variantes da etapa de
classificação: (i)~a \textbf{executada} (Seção~\ref{sec:score-exec}), em que
o \textit{proxy} de jargão é a contagem de sinais morfológicos greco-latinos;
e (ii)~a \textbf{proposta} (Seção~\ref{sec:score-prop}), em que o
\textit{proxy} é a quantidade de expressões encontradas no DeCS. Os
resultados reportados na Seção~\ref{sec:resultados} correspondem à variante
executada. A variante com DeCS está em espera porque o XML completo foi
solicitado à BIREME/BVS, mas ainda não foi liberado até o fechamento desta
entrega; a consulta termo a termo é inviável para 17.182 registros.

\subsection{Fontes de dados}

Utilizamos duas tabelas oficiais do DATASUS, descritas na
Tabela~\ref{tab:fontes}. O relacionamento entre elas não é explorado aqui:
CID-10 e SIGTAP são tratados como duas fontes de \textit{texto}
independentes, pois o objetivo é medir a linguagem.

\begin{table}[ht]
\centering
\caption{Fontes de dados utilizadas (DATASUS).}
\label{tab:fontes}
\begin{tabular}{lll}
\toprule
\textbf{Fonte} & \textbf{Conteúdo} & \textbf{Registros} \\
\midrule
CID-10 (V2008)         & diagnósticos (código + nome) & 12.198 \\
SIGTAP (comp. 05/2026) & procedimentos (código + nome) & 4.984 \\
\bottomrule
\end{tabular}
\end{table}

\subsection{Coleta, integração e limpeza}

A CID-10~\cite{datasus_cid10} e o SIGTAP~\cite{datasus_sigtap} foram obtidos
em arquivos de \textit{largura fixa} (sem separador), lidos por posição
conforme seus respectivos \textit{lay-outs}. As duas fontes foram
concatenadas em uma única tabela, preservando a origem de cada registro.
A limpeza tratou valores ausentes e duplicidades. A detecção de duplicidades
é feita por \textbf{descrição}, e não pela combinação \texttt{(codigo,
descricao)}: como os códigos da CID-10 e do SIGTAP são únicos por
construção, a verificação combinada sempre devolveria zero; o que importa
para o nosso objetivo --- medir linguagem --- é a repetição da descrição.
Após a limpeza, o \textit{dataset} contém \textbf{17.182 registros}.

\subsection{Extração de atributos}

Para cada descrição são derivados seis atributos textuais:
\texttt{qtd\_palavras}, \texttt{qtd\_caracteres}, \texttt{qtd\_palavras\_longas}
(palavras com $\geq 12$ caracteres), \texttt{qtd\_sinais\_morfologicos},
\texttt{tem\_sigla} e \texttt{tem\_abreviacao}. Os atributos
\texttt{qtd\_palavras\_longas} e \texttt{qtd\_sinais\_morfologicos} são
mantidos \emph{separadamente} para evitar confundir tamanho com jargão:
``tratamento'' é uma palavra longa de uso comum, enquanto ``gastrite''
é curta, mas carrega um sinal morfológico médico explícito.

\paragraph{Sinais morfológicos.} Adotamos uma lista de morfemas
greco-latinos típicos da linguagem médica~\cite{rezende2004} (por exemplo,
\textit{-ite}, \textit{-ectomia}, \textit{-scopia}, \textit{cardio-},
\textit{nefro-}). Uma palavra é contada como sinal morfológico se contiver
ao menos um desses morfemas e não estiver em uma lista de palavras
funcionais comuns.

\paragraph{Detecção de siglas.} A detecção de siglas usa um
\emph{vocabulário explícito e auditável} de siglas comuns no SUS (SUS, AIDS,
DPOC, AVC, IAM, ECG, UTI etc.), em vez de uma expressão regular sobre
palavras em caixa alta. O SIGTAP é distribuído majoritariamente em
\textsc{caixa alta}; uma regex baseada apenas em maiúsculas confundiria
palavras comuns com siglas.

\subsection{Correspondência com o DeCS (variante proposta)}
\label{sec:decs}

O atributo principal da variante proposta é \texttt{qtd\_termos\_decs}: a
quantidade de termos ou expressões completas encontrados no vocabulário
controlado DeCS~\cite{decs}. O DeCS é usado como tabela de consulta; suas
entradas não são adicionadas como novas linhas ao \textit{dataset}. A busca
normaliza as descrições e prioriza a expressão mais longa encontrada,
evitando contar ao mesmo tempo uma expressão e seus componentes. Esta
variante depende da liberação do XML completo pela BIREME/BVS.

\subsection{Escore e faixas exploratórias}

\subsubsection{Variante executada}
\label{sec:score-exec}

Os atributos numéricos são normalizados para o intervalo $[0,1]$ e
combinados conforme a Equação~\ref{eq:score-exec}:
\begin{equation}
\label{eq:score-exec}
\textit{score} = 0{,}40\,M + 0{,}20\,P + 0{,}15\,P_L + 0{,}10\,C +
0{,}10\,S + 0{,}05\,A,
\end{equation}

\noindent onde $M$ é \texttt{qtd\_sinais\_morfologicos}, $P$ é
\texttt{qtd\_palavras}, $P_L$ é \texttt{qtd\_palavras\_longas}, $C$ é
\texttt{qtd\_caracteres}, $S$ é \texttt{tem\_sigla} e $A$ é
\texttt{tem\_abreviacao}. A maior ponderação de $M$ é uma hipótese
operacional motivada por evidências de que jargão melhora a explicação da
dificuldade em textos médicos~\cite{jiang2024medreadme}.

\subsubsection{Variante proposta (com DeCS)}
\label{sec:score-prop}

Quando o XML do DeCS for liberado, o \textit{proxy} de jargão $M$ é
substituído por $T_{DeCS}$ (Eq.~\ref{eq:score-prop}), assumindo a maior
ponderação:
\begin{equation}
\label{eq:score-prop}
\textit{score} = 0{,}50\,T_{DeCS} + 0{,}15\,M + 0{,}20\,P +
0{,}10\,S + 0{,}05\,A.
\end{equation}

Em ambas as variantes, os pesos exatos não são tratados como escala
validada. O escore contínuo é dividido em tercis para formar faixas
exploratórias baixa, média e alta. Como esse procedimento equilibra as
classes por construção, ele não mede a prevalência natural de complexidade
e não substitui rótulos humanos.

\subsection{Validação por anotação humana}
\label{sec:validacao}

Será exportada uma amostra estratificada de 150 descrições, sendo 75 de cada
fonte e distribuídas pelas três faixas exploratórias (25 por faixa por fonte).
O tamanho foi escolhido para equilibrar viabilidade da rodada com sensibilidade
adequada para estimar kappa de classes ordinais; uma rodada-piloto de 20 itens
definirá o guia de anotação. A amostra não recebe o escore previamente, para
evitar viés de ancoragem. Dois avaliadores leigos --- representando o público
de referência, adultos brasileiros sem formação em saúde --- rotularão
independentemente cada descrição. A concordância será medida pelo
kappa ponderado de Cohen, adequado a classes ordinais, e interpretada pelos
patamares de Landis e Koch~\cite{landis1977}. O consenso humano será então
comparado às faixas heurísticas.

\section{Resultados}
\label{sec:resultados}

\subsection{Qualidade e composição da base}

O \textit{dataset} contém 17.182 registros: 12.198 (71{,}0\%) da CID-10 e
4.984 (29{,}0\%) do SIGTAP. Não foram identificados códigos vazios nem
descrições vazias. Há descrições idênticas com códigos diferentes,
sobretudo no SIGTAP (formas organizacionais distintas para o mesmo
procedimento); para esta entrega, a duplicidade semântica por descrição é
explicitamente mensurada e tratada. Em uma futura modelagem supervisionada,
descrições idênticas deverão permanecer no mesmo grupo de treino ou teste
para evitar vazamento de informação.

\subsection{Estatísticas descritivas preliminares}
\label{sec:descritivos}

A Tabela~\ref{tab:descritivos} apresenta atributos calculados na variante
executada do \textit{pipeline}. A CID-10 apresenta maior extensão média, em
palavras e caracteres. O SIGTAP apresenta, em média, mais palavras longas e
sinais morfológicos médicos.

\begin{table}[ht]
\centering
\caption{Estatísticas descritivas dos atributos textuais por fonte.}
\label{tab:descritivos}
\begin{tabular}{lrrrrrr}
\toprule
& \multicolumn{2}{c}{Palavras} & \multicolumn{2}{c}{Caracteres} &
Palavras longas & Sinais morfológicos \\
\cmidrule(lr){2-3}\cmidrule(lr){4-5}\cmidrule(lr){6-6}\cmidrule(lr){7-7}
Fonte & Média & Mediana & Média & Mediana & Média & Média \\
\midrule
CID-10 & 7,66 & 7,00 & 53,47 & 51,00 & 0,59 & 0,24 \\
SIGTAP & 6,75 & 6,00 & 49,43 & 45,50 & 0,66 & 0,31 \\
\bottomrule
\end{tabular}
\end{table}

Esses resultados impedem uma conclusão simplista de que procedimentos são
automaticamente mais complexos por serem procedimentos. A CID-10 possui maior
extensão média; o SIGTAP possui valores ligeiramente maiores em dois sinais
auxiliares de especialização. Somente a correspondência com o DeCS e a
validação humana permitirão discutir complexidade técnico-linguística de
forma mais fundamentada.

A maior descrição SIGTAP possui 249 caracteres e 40 palavras, enquanto a
maior descrição CID-10 possui 100 caracteres e 21 palavras. Isso sugere
maior heterogeneidade de extensão no SIGTAP, mas máximos não substituem
análise do conjunto.

\subsection{Resultados dependentes do DeCS e da anotação humana}

Após o recebimento da versão licenciada do DeCS, esta subseção deverá
incluir: percentual de descrições com pelo menos uma expressão encontrada;
distribuição de \texttt{qtd\_termos\_decs}; termos mais frequentes;
distribuição do escore; e distribuição das faixas por fonte. Após a anotação
humana, devem ser reportados kappa ponderado, concordância simples e matriz
de confusão entre o consenso humano e a faixa heurística.

Uma árvore de decisão pode ser utilizada para visualizar quais atributos
reproduzem a regra. Entretanto, não se deve interpretá-la como classificador
preditivo enquanto o alvo for produzido pelos mesmos atributos utilizados
como entrada.

\section{Discussão e Limitações}

O \textit{dataset} construído permite caracterizar, de forma reproduzível,
atributos associados à complexidade técnico-linguística das descrições do
SUS. Reconhecemos limitações: (i)~a complexidade é uma \textit{proxy}
heurística, não uma medida direta de compreensão por pacientes;
(ii)~a presença no DeCS indica terminologia de saúde, mas não garante
dificuldade para todo leitor; (iii)~o alvo deriva dos atributos, o que torna
uma árvore apenas interpretativa; e (iv)~o equilíbrio entre classes é por
construção dos tercis.

\textbf{Dificuldades enfrentadas.} A construção exigiu superar obstáculos
práticos: a instabilidade dos portais do DATASUS; a descontinuação do
suporte a FTP nos navegadores modernos (contornada via Explorador de
Arquivos); o formato de largura fixa do SIGTAP; a atualização mensal por
competência; e a tradução oficial da CID-10 congelada em 2008, porém
vigente até 2027 (Nota Técnica MS 91/2024). A integração do DeCS, originalmente
prevista para esta entrega, foi adiada devido ao tempo de processamento da
solicitação do XML completo à BIREME/BVS.

\section{Reprodutibilidade e Ameaças à Validade}

A execução final registrará versões das fontes, competência do SIGTAP, edição
do DeCS, critérios de limpeza, lista de siglas, pesos do escore e semente da
amostragem. A validade de construto é limitada porque presença no DeCS não
prova dificuldade para todo leitor. A validade externa restringe-se ao
português brasileiro, às versões analisadas e a descrições administrativas
isoladas. A validação humana proposta reduz, mas não elimina, essas
limitações.

\section{Considerações Finais}

Este trabalho construiu uma base de descrições e um \textit{pipeline}
reproduzível para análise exploratória da complexidade técnico-linguística
de diagnósticos e procedimentos do SUS, com duas variantes explícitas:
uma executada (com proxy morfológico) e uma proposta (com DeCS). Os
resultados descritivos mostram que as diferenças entre CID-10 e SIGTAP não
podem ser resumidas apenas pela extensão textual. Como trabalhos futuros,
pretende-se: executar a correspondência com o DeCS assim que o XML for
liberado e validar as faixas com múltiplos anotadores; comparar a
complexidade linguística com a complexidade assistencial do SIGTAP; e
integrar o classificador a uma ferramenta de simplificação de linguagem
médica.

\bibliographystyle{sbc}
\bibliography{references}

\end{document}
